\documentclass[aspectratio=169, xcolor=dvipsnames]{beamer}

% --- Theme Setup ---
\usetheme{Madrid}
\usecolortheme{default}

% --- Packages ---
\usepackage{graphicx}
\usepackage{xcolor}
\usepackage{booktabs}
\usepackage{hyperref}
\usepackage{adjustbox}
\usepackage{amssymb}
\usepackage{amsmath}
\usepackage{listings}
\usepackage{caption}
\usepackage{tikz}

\lstset{
  basicstyle=\ttfamily\small,
  keywordstyle=\color{blue}\bfseries,
  stringstyle=\color{red},
  showstringspaces=false,
  breaklines=true
}

% --- Title Information ---
\title{Module 26}
\subtitle{Relational Database Design/6: Normal Forms}
\author{Milav Dabgar}
\institute{www.milav.in}
\date{\today}

\begin{document}

% Title Slide
\begin{frame}
    \titlepage
\end{frame}

\begin{frame}{Week Recap}
    \begin{itemize}[<+->]
        \item Identified the features of good relational design
        \item Familiarized with the First Normal Form
        \item Introduced the notion and the theory of functional dependencies
        \item Discussed issues in `good' design in the context of functional dependencies
        \item Studied Algorithms for Properties of Functional Dependencies
        \item Understood the Characterization for and Determination of Lossless Join and Determination of Dependency Preservation
    \end{itemize}
\end{frame}

\begin{frame}{Module Objectives}
    \begin{block}{Goals}
        \begin{itemize}[<+->]
            \item To Understand the Normal Forms and their Importance in Relational Design
        \end{itemize}
    \end{block}
\end{frame}

\begin{frame}{Module Outline}
    \begin{itemize}[<+->]
        \item Normal Forms
        \begin{itemize}
            \item[$\triangleright$] 1NF
            \item[$\triangleright$] 2NF
            \item[$\triangleright$] 3NF
        \end{itemize}
    \end{itemize}
\end{frame}

\section{Normal Forms}
\begin{frame}
  \vfill
  \centering
  \begin{beamercolorbox}[sep=8pt,center,shadow=true,rounded=true]{title}
    \usebeamerfont{title}Normal Forms\par%
  \end{beamercolorbox}
  \vfill
\end{frame}

\begin{frame}{Normalization or Schema Refinement}
    \begin{itemize}[<+->]
        \item Normalization or Schema Refinement is a technique of organizing the data in the database
        \item A systematic approach of decomposing tables to eliminate data redundancy and undesirable characteristics
        \begin{itemize}
            \item[$\triangleright$] Insertion Anomaly
            \item[$\triangleright$] Update Anomaly
            \item[$\triangleright$] Deletion Anomaly
        \end{itemize}
        \item Most common technique for the Schema Refinement is decomposition.
        \item \textbf{Goal of Normalization:} Eliminate Redundancy
        \item Redundancy refers to repetition of same data or duplicate copies of same data stored in different locations
        \item Normalization is used for mainly two purpose:
        \begin{itemize}
            \item[$\triangleright$] Eliminating redundant (useless) data
            \item[$\triangleright$] Ensuring data dependencies make sense, that is, data is logically stored
        \end{itemize}
    \end{itemize}
\end{frame}

\begin{frame}{Anomalies}
    \textbf{Employees' Skills}
    \begin{itemize}[<+->]
        \item a) \textbf{Update Anomaly}: Employee 519 is shown as having different addresses on different records
        \item b) \textbf{Insertion Anomaly}: Until the new faculty member, Dr. Newsome, is assigned to teach at least one course, his details cannot be recorded in Faculty and Their Courses
        \item c) \textbf{Deletion Anomaly}: All information about Dr. Giddens is lost if he temporarily ceases to be assigned to any courses.
    \end{itemize}
    \vspace{0.5em}
    \textbf{Resolution:} Decompose the Schema
    \begin{itemize}[<+(3)->]
        \item a) Update: \texttt{(ID, Address)}, \texttt{(ID, Skill)}
        \item b) Insert: \texttt{(ID, Name, Hire Date)}, \texttt{(ID, Code)}
        \item c) Delete: \texttt{(ID, Name, Hire Date)}, \texttt{(ID, Code)}
    \end{itemize}
\end{frame}

\begin{frame}{Desirable Properties of Decomposition}
    \begin{itemize}[<+->]
        \item \textbf{Lossless Join Decomposition Property}
        \begin{itemize}
            \item[$\triangleright$] It should be possible to reconstruct the original table
        \end{itemize}
        \item \textbf{Dependency Preserving Property}
        \begin{itemize}
            \item[$\triangleright$] No functional dependency (or other constraints should get violated)
        \end{itemize}
    \end{itemize}
\end{frame}

\begin{frame}{Normalization and Normal Forms}
    \begin{itemize}[<+->]
        \item A normal form specifies a set of conditions that the relational schema must satisfy in terms of its constraints - they offer varied levels of guarantee for the design
        \item Normalization rules are divided into various normal forms. Most common normal forms are:
        \begin{itemize}
            \item[$\triangleright$] First Normal Form (1NF)
            \item[$\triangleright$] Second Normal Form (2NF)
            \item[$\triangleright$] Third Normal Form (3NF)
        \end{itemize}
        \item Informally, a relational database relation is often described as `normalized' if it meets third normal form. Most 3NF relations are free of insertion, update, and deletion anomalies.
    \end{itemize}
\end{frame}

\begin{frame}{Normalization and Normal Forms (Continued)}
    Additional Normal Forms:
    \begin{itemize}[<+->]
        \item Elementary Key Normal Form (EKNF)
        \item Boyce-codd Normal Form (BCNF)
        \item Multivalued Dependencies And Fourth Normal Form (4NF)
        \item Essential Tuple Normal Form (ETNF)
        \item Join Dependencies and Fifth Normal Form (5NF)
        \item Sixth Normal Form (6NF)
        \item Domain/Key Normal Form (DKNF)
    \end{itemize}
\end{frame}

\section{1NF}
\begin{frame}
  \vfill
  \centering
  \begin{beamercolorbox}[sep=8pt,center,shadow=true,rounded=true]{title}
    \usebeamerfont{title}1NF\par%
  \end{beamercolorbox}
  \vfill
\end{frame}

\begin{frame}{1NF: First Normal Form}
    \begin{block}{Definition}
        \begin{itemize}[<+->]
            \item A relation is in First Normal Form if and only if all underlying domains contain atomic values only (doesn't have multivalued attributes (MVA))
            \item Example: \texttt{STUDENT(Sid, Sname, Cname)}
        \end{itemize}
    \end{block}
    \vspace{0.5em}
    \begin{columns}
        \column{0.45\textwidth}
        \textbf{Not in 1NF (MVA)}
        \begin{table}
        \resizebox{\linewidth}{!}{
        \begin{tabular}{lll}
        \toprule
        \textbf{SID} & \textbf{Sname} & \textbf{Cname} \\
        \midrule
        S1 & A & C, C++ \\
        S2 & B & C++, DB \\
        S3 & A & DB \\
        \bottomrule
        \end{tabular}
        }
        \end{table}
        \column{0.45\textwidth}
        \textbf{In 1NF (No MVA)}
        \begin{table}
        \resizebox{\linewidth}{!}{
        \begin{tabular}{lll}
        \toprule
        \textbf{SID} & \textbf{Sname} & \textbf{Cname} \\
        \midrule
        S1 & A & C \\
        S1 & A & C++ \\
        S2 & B & C++ \\
        S2 & B & DB \\
        S3 & A & DB \\
        \bottomrule
        \end{tabular}
        }
        \end{table}
    \end{columns}
\end{frame}

\begin{frame}{1NF (2): Possible Redundancy}
    \begin{itemize}
        \item Example: \texttt{Supplier(SID, Status, City, PID, Qty)}
    \end{itemize}
    \begin{table}
    \centering
    \begin{tabular}{lllll}
    \toprule
    \textbf{SID} & \textbf{Status} & \textbf{City} & \textbf{PID} & \textbf{Qty} \\
    \midrule
    S1 & 30 & Delhi & P1 & 100 \\
    S1 & 30 & Delhi & P2 & 125 \\
    S1 & 30 & Delhi & P3 & 200 \\
    S1 & 30 & Delhi & P4 & 130 \\
    S2 & 10 & Karnal & P1 & 115 \\
    S2 & 10 & Karnal & P2 & 250 \\
    S3 & 40 & Rohtak & P1 & 245 \\
    S4 & 30 & Delhi & P4 & 300 \\
    S4 & 30 & Delhi & P5 & 315 \\
    \bottomrule
    \multicolumn{5}{l}{\footnotesize Key (SID, PID)}
    \end{tabular}
    \end{table}
\end{frame}

\begin{frame}{1NF (2): Drawbacks}
    \begin{alertblock}{Drawbacks of 1NF}
        \begin{itemize}[<+->]
            \item \textbf{Deletion Anomaly:} If we delete \texttt{<S3, 40, Rohtak, P1, 245>}, then we lose the information that S3 lives in Rohtak.
            \item \textbf{Insertion Anomaly:} We cannot insert a Supplier S5 located in Karnal, until S5 supplies at least one part.
            \item \textbf{Update Anomaly:} If Supplier S1 moves from Delhi to Kanpur, then it is difficult to update all the tuples having SID as S1 and City as Delhi.
        \end{itemize}
    \end{alertblock}
    \vspace{0.5em}
    \textit{Normalization is a method to reduce redundancy. However, sometimes 1NF increases redundancy.}
\end{frame}

\begin{frame}{1NF (3): Possible Redundancy}
    \textbf{When LHS is not a Superkey:}
    \begin{itemize}[<+->]
        \item Let $X \rightarrow Y$ be a non trivial FD over $R$ with $X$ is not a superkey of $R$, then redundancy exist between $X$ and $Y$ attribute set.
        \item Hence in order to identify the redundancy, we need not to look at the actual data, it can be identified by given functional dependency.
        \item Example: $X \rightarrow Y$ and $X$ is not a Candidate Key
        \begin{itemize}
            \item[$\triangleright$] $\Rightarrow X$ can duplicate
            \item[$\triangleright$] $\Rightarrow$ Corresponding $Y$ value would duplicate also.
        \end{itemize}
    \end{itemize}
    \vspace{0.5em}
    \textbf{When LHS is a Superkey:}
    \begin{itemize}[<+(4)->]
        \item If $X \rightarrow Y$ is a non trivial FD over $R$ with $X$ is a superkey of $R$, then redundancy does not exist between $X$ and $Y$ attribute set.
        \item Example: $X \rightarrow Y$ and $X$ is a Candidate Key $\Rightarrow X$ cannot duplicate
        \begin{itemize}
            \item[$\triangleright$] $\Rightarrow$ Corresponding $Y$ value may or may not duplicate.
        \end{itemize}
    \end{itemize}
\end{frame}

\section{2NF}
\begin{frame}
  \vfill
  \centering
  \begin{beamercolorbox}[sep=8pt,center,shadow=true,rounded=true]{title}
    \usebeamerfont{title}2NF\par%
  \end{beamercolorbox}
  \vfill
\end{frame}

\begin{frame}{2NF: Second Normal Form}
    \begin{block}{Definition of 2NF}
        \begin{itemize}[<+->]
            \item Relation $R$ is in Second Normal Form (2NF) only iff:
            \begin{itemize}
                \item[$\triangleright$] $R$ is in 1NF and
                \item[$\triangleright$] $R$ contains no Partial Dependency
            \end{itemize}
        \end{itemize}
    \end{block}
    \vspace{0.5em}
    \textbf{Partial Dependency:}
    \begin{itemize}[<+(2)->]
        \item Let $R$ be a relational Schema and $X, Y, A$ be the attribute sets over $R$ where $X$: Any Candidate Key, $Y$: Proper Subset of Candidate Key, and $A$: Non Prime Attribute
        \item If $Y \rightarrow A$ exists in $R$, then $R$ is not in 2NF.
    \end{itemize}
    $(Y \rightarrow A)$ is a Partial dependency only if:
    \begin{itemize}[<+(4)->]
        \item $Y$: Proper subset of Candidate Key
        \item $A$: Non Prime Attribute
        \item \textit{A prime attribute of a relation is an attribute that is a part of a candidate key of the relation}
    \end{itemize}
\end{frame}

\begin{frame}{2NF (2): Possible Redundancy}
    \begin{itemize}
        \item \texttt{STUDENT(Sid, Sname, Cname)} (already in 1NF)
        \item Redundancy? Yes: Sname
        \item Anomaly? Yes
    \end{itemize}
    \begin{table}
    \centering
    \begin{tabular}{lll}
    \toprule
    \textbf{SID} & \textbf{Sname} & \textbf{Cname} \\
    \midrule
    S1 & A & C \\
    S1 & A & C++ \\
    S2 & B & C++ \\
    S2 & B & DB \\
    S3 & A & DB \\
    \bottomrule
    \multicolumn{3}{l}{\footnotesize (SID, Cname): Primary Key}
    \end{tabular}
    \end{table}
    \textbf{Functional Dependencies:} $\{SID, Cname\} \rightarrow Sname$, $SID \rightarrow Sname$ \\
    \textbf{Partial Dependencies:} $SID \rightarrow Sname$ (as SID is a Proper Subset of Candidate Key $\{SID, Cname\}$)
\end{frame}

\begin{frame}{2NF (3): Possible Redundancy}
    \begin{itemize}
        \item \texttt{Supplier(SID, Status, City, PID, Qty)}
    \end{itemize}
    \begin{table}
    \centering
    \resizebox{\linewidth}{!}{
    \begin{tabular}{lllll}
    \toprule
    \textbf{SID} & \textbf{Status} & \textbf{City} & \textbf{PID} & \textbf{Qty} \\
    \midrule
    S1 & 30 & Delhi & P1 & 100 \\
    S1 & 30 & Delhi & P2 & 125 \\
    S1 & 30 & Delhi & P3 & 200 \\
    S1 & 30 & Delhi & P4 & 130 \\
    S2 & 10 & Karnal & P1 & 115 \\
    S2 & 10 & Karnal & P2 & 250 \\
    S3 & 40 & Rohtak & P1 & 245 \\
    S4 & 30 & Delhi & P4 & 300 \\
    S4 & 30 & Delhi & P5 & 315 \\
    \bottomrule
    \multicolumn{5}{l}{\footnotesize Key (SID, PID)}
    \end{tabular}
    }
    \end{table}
    \textbf{Partial Dependencies:} \texttt{SID $\rightarrow$ Status}, \texttt{SID $\rightarrow$ City}
\end{frame}

\begin{frame}{2NF (3): Drawbacks}
    \begin{alertblock}{Drawbacks}
        \begin{itemize}[<+->]
            \item \textbf{Deletion Anomaly}: If we delete a tuple in Sup City, then we not only loose the information about a supplier, but also loose the status value of a particular city.
            \item \textbf{Insertion Anomaly}: We cannot insert a City and its status until a supplier supplies at least one part.
            \item \textbf{Update Anomaly}: If the status value for a city is changed, then we will face the problem of searching every tuple for that city.
        \end{itemize}
    \end{alertblock}
\end{frame}

\section{3NF}
\begin{frame}
  \vfill
  \centering
  \begin{beamercolorbox}[sep=8pt,center,shadow=true,rounded=true]{title}
    \usebeamerfont{title}3NF\par%
  \end{beamercolorbox}
  \vfill
\end{frame}

\begin{frame}{3NF: Third Normal Form}
    Let $R$ be the relational schema.
    \begin{block}{Definitions}
        \begin{itemize}[<+->]
            \item \textbf{[E. F. Codd, 1971]} $R$ is in 3NF only if:
            \begin{itemize}
                \item[$\triangleright$] $R$ should be in 2NF
                \item[$\triangleright$] $R$ should not contain transitive dependencies (OR, Every non-prime attribute of $R$ is non-transitively dependent on every key of $R$)
            \end{itemize}
            \item \textbf{[Carlo Zaniolo, 1982]} Alternately, $R$ is in 3NF iff for each of its functional dependencies $X \rightarrow A$, at least one of the following conditions holds:
            \begin{itemize}
                \item[$\triangleright$] $X$ contains $A$ (that is, $A$ is a subset of $X$, meaning $X \rightarrow A$ is trivial functional dependency), or
                \item[$\triangleright$] $X$ is a superkey, or
                \item[$\triangleright$] Every element of $A - X$, the set difference between $A$ and $X$, is a prime attribute (i.e., each attribute in $A - X$ is contained in some candidate key)
            \end{itemize}
            \item \textbf{[Simple Statement]} A relational schema $R$ is in 3NF if for every FD $X \rightarrow A$ associated with $R$ either:
            \begin{itemize}
                \item[$\triangleright$] $A \subseteq X$ (that is, the FD is trivial) or
                \item[$\triangleright$] $X$ is a superkey of $R$ or
                \item[$\triangleright$] $A$ is part of some candidate key (not just superkey!)
            \end{itemize}
            \item A relation in 3NF is naturally in 2NF
        \end{itemize}
    \end{block}
\end{frame}

\begin{frame}{3NF (2): Transitive Dependency}
    \begin{itemize}[<+->]
        \item A transitive dependency is a functional dependency which holds by virtue of transitivity. A transitive dependency can occur only in a relation that has three or more attributes.
        \item Let $A$, $B$, and $C$ designate three distinct attributes (or distinct collections of attributes) in the relation. Suppose all three of the following conditions hold:
        \begin{enumerate}
            \item $A \rightarrow B$
            \item It is not the case that $B \rightarrow A$
            \item $B \rightarrow C$
        \end{enumerate}
        \item Then the functional dependency $A \rightarrow C$ (which follows from 1 and 3 by the axiom of transitivity) is a transitive dependency
    \end{itemize}
\end{frame}

\begin{frame}{3NF (3): Transitive Dependency}
    \textbf{Example of transitive dependency}
    \begin{itemize}[<+->]
        \item The functional dependency \texttt{\{Book\} $\rightarrow$ \{Author Nationality\}} applies; that is, if we know the book, we know the author's nationality. Furthermore:
        \begin{itemize}
            \item[$\triangleright$] \texttt{\{Book\} $\rightarrow$ \{Author\}}
            \item[$\triangleright$] \texttt{\{Author\}} does not $\rightarrow$ \texttt{\{Book\}}
            \item[$\triangleright$] \texttt{\{Author\} $\rightarrow$ \{Author Nationality\}}
        \end{itemize}
        \item Therefore \texttt{\{Book\} $\rightarrow$ \{Author Nationality\}} is a transitive dependency.
        \item Transitive dependency occurred because a non-key attribute (Author) was determining another non-key attribute (Author Nationality).
    \end{itemize}
    \begin{table}
    \centering
    \resizebox{0.9\linewidth}{!}{
    \begin{tabular}{llll}
    \toprule
    \textbf{Book} & \textbf{Genre} & \textbf{Author} & \textbf{Author Nationality} \\
    \midrule
    Twenty Thousand Leagues... & Science Fiction & Jules Verne & French \\
    Journey to the Center... & Science Fiction & Jules Verne & French \\
    Leaves of Grass & Poetry & Walt Whitman & American \\
    Anna Karenina & Literary Fiction & Leo Tolstoy & Russian \\
    A Confession & Religious Autobiography & Leo Tolstoy & Russian \\
    \bottomrule
    \end{tabular}
    }
    \end{table}
\end{frame}

\begin{frame}{3NF (4): Example}
    \textbf{Example:} \texttt{Sup\_City(SID, Status, City)} (already in 2NF)
    \begin{itemize}[<+->]
        \item \textbf{Functional Dependencies:} \texttt{SID $\rightarrow$ Status}, \texttt{SID $\rightarrow$ City}, \texttt{City $\rightarrow$ Status}
        \item \textbf{Transitive Dependency:} \texttt{SID $\rightarrow$ Status} \{As \texttt{SID $\rightarrow$ City} and \texttt{City $\rightarrow$ Status}\}
    \end{itemize}
    \begin{table}
    \centering
    \begin{tabular}{lll}
    \toprule
    \textbf{SID} & \textbf{City} & \textbf{Status} \\
    \midrule
    S1 & Delhi & 30 \\
    S2 & Karnal & 10 \\
    S3 & Rohtak & 40 \\
    S4 & Delhi & 30 \\
    \bottomrule
    \multicolumn{3}{l}{\footnotesize SID: Primary Key}
    \end{tabular}
    \end{table}
\end{frame}

\begin{frame}{Module Summary}
    \begin{block}{Summary}
        \begin{itemize}[<+->]
            \item To Understand the Normal Forms and their Importance in Relational Design
        \end{itemize}
    \end{block}
    \vspace{2em}
    \begin{center}
    \small \textit{Slides used in this presentation are borrowed from \url{http://db-book.com/} with kind permission of the authors.} \\
    \small \textit{Edited and new slides are marked with 'PPD'.}
    \end{center}
\end{frame}

\end{document}
